%!TEX root = thesis.tex

\chapter{Evaluierung}
\label{chapter-evaluation}

In der Evaluierung wird das Ergebnis dieser Arbeit bewertet. Eine praktische Evaluation eines neuen Algorithmus kann zum Beispiel durch eine Implementierung geschehen. Je nach Thema der Arbeit kann sich natürlich auch die gesamte Arbeit eher im praktischen Bereich mit einer Implementierung beschäftigen. In diesem Fall gilt es am Ende der Arbeit insbesondere die Implementierung selber zu evaluieren. Wesentliche Fragen dabei können sein:
\begin{compactitem}[--]
  \item Was funktioniert jetzt besser als vor meiner Arbeit?
  \item Wie kann das praktisch eingesetzt werden?
  \item Was sagen potenzielle Anwender zu meiner Lösung?
\end{compactitem}

\section{Ergebnisse und Auswertung}
In diesem Kapitel werden die Ergebnisse der durchgeführten
Experimente dargestellt und analysiert. Ziel ist es,
die durch verschiedene Large Language Models generierten
Softwaresysteme systematisch zu vergleichen und im Hinblick
auf die definierten Evaluationskriterien zu bewerten.

Die Ergebnisse basieren auf insgesamt neun Experimenten,
die sich aus der Kombination von drei Prompt-Varianten
und drei unterschiedlichen LLM-Modellen ergeben.

Im Folgenden werden zunächst die Ergebnisse der einzelnen
Experimente vorgestellt. Anschließend erfolgt ein
vergleichender Überblick sowie eine Diskussion der
gewonnenen Erkenntnisse.

\section{Ergebnisse der einzelnen Experimente}
\subsection{Experiment 1: Prompt 1 mit GPT-5.2}

Das generierte System basiert auf dem Django-Framework
und folgt grundsätzlich der Model-View-Template-Architektur.

Die grundlegenden Funktionen, wie die Verwaltung von
Kundendaten und Buchungen, wurden umgesetzt. Allerdings
zeigten sich Einschränkungen bei der Umsetzung komplexerer
Anforderungen.

Die Codequalität ist insgesamt als mittelmäßig zu bewerten,
da zwar eine grundlegende Struktur vorhanden ist, jedoch
Inkonsistenzen in der Implementierung auftreten.

Die Benutzeroberfläche ist funktional, jedoch einfach
gehalten und weist nur eine begrenzte Benutzerfreundlichkeit
auf.

\subsection{Experiment 1: GPT-5.2 mit Basis-Prompt}

Im ersten Experiment wurde das Buchungssystem
mithilfe des Modells GPT-5.2 generiert. Das resultierende
System wurde als vollständige Django-Webanwendung
implementiert und anschließend hinsichtlich der definierten
Evaluationskriterien analysiert.

\textbf{Funktionalität und Vollständigkeit:}
Die Analyse zeigt, dass das generierte System sämtliche
zentralen Anforderungen des ursprünglichen Legacy-Systems
umsetzt. Dazu zählen insbesondere die Verwaltung von
Kundendaten, die Buchungsverwaltung, die
Rechnungserstellung sowie eine Kalenderübersicht.
Die vollständige Umsetzung der Spezifikation wird auch
durch die vorhandene Dokumentation bestätigt.

\textbf{Architektur:}
Das System folgt konsequent der Model-View-Template-
Architektur des Django-Frameworks. Die klare Trennung
zwischen Datenmodellen, Logik und Präsentationsschicht
weist auf eine korrekte Umsetzung grundlegender
Architekturprinzipien hin.

\textbf{Codequalität:}
Die generierte Codebasis ist strukturiert und modular
aufgebaut. Die Verwendung mehrerer Komponenten, wie
Models, Views und Templates, ermöglicht eine gute
Nachvollziehbarkeit und Erweiterbarkeit des Systems.
Dennoch sind kleinere Optimierungspotenziale im Hinblick
auf Code-Konsistenz und Best Practices erkennbar.

\textbf{Benutzeroberfläche:}
Die grafische Benutzeroberfläche basiert auf Bootstrap
und bietet eine übersichtliche und funktionale
Interaktion. Das Dashboard sowie die Formularstrukturen
ermöglichen eine intuitive Nutzung, auch wenn die
Gestaltung eher funktional als innovativ ist.

\textbf{Performance und Skalierbarkeit:}
Das System verwendet standardmäßig eine SQLite-Datenbank,
was für Entwicklungszwecke ausreichend ist, jedoch
Einschränkungen hinsichtlich Skalierbarkeit und
Performance aufweist.

\textbf{Sicherheit und Wartbarkeit:}
Grundlegende Sicherheitsmechanismen, wie Login-Schutz
und CSRF-Schutz, sind implementiert. Die modulare
Struktur des Systems unterstützt zudem die Wartbarkeit
und zukünftige Erweiterungen.

\textbf{Zusammenfassung der Bewertung:}

\begin{table}[h]
\centering
\begin{tabular}{lc}
\hline
\textbf{Kriterium} & \textbf{Score (0–5)} \\
\hline
Funktionalität & 5 \\
Vollständigkeit & 5 \\
Codequalität & 4 \\
Architektur & 5 \\
Benutzeroberfläche & 4 \\
Performance & 3 \\
Sicherheit & 4 \\
Wartbarkeit & 4 \\
\hline
\end{tabular}
\caption{Bewertung von Experiment 1 (GPT-5.2)}
\end{table}

Insgesamt zeigt das erste Experiment, dass GPT-5.2 in
der Lage ist, ein nahezu vollständig funktionsfähiges
Softwaresystem zu generieren. Die Qualität der Ergebnisse
ist insbesondere im Bereich der Funktionalität und
Architektur sehr hoch, während kleinere Einschränkungen
in den Bereichen Performance und Codeoptimierung bestehen.

\section{Implementierungen}

Wenn Implementierungen umfangreich beschrieben werden, ist darauf zu achten, den richtigen Mittelweg zwischen einer zu detaillierten und zu oberflächlichen Beschreibung zu finden. Eine Beschreibung aller Details der Implementierung ist in der Regel zu detailliert, da die primäre Zielgruppe einer Abschlussarbeit sich nicht im Detail in den geschriebenen Quelltext einarbeiten will. Die Beschreibung sollte aber durchaus alle wesentlichen Konzepte der Implementierung enthalten. Gerade bei einer Abschlussarbeit am Institut für Softwaretechnik und Programmiersprachen lohnt es sich, auf die eingesetzten Techniken und Programmiersprachen einzugehen. Ich würde in einer solchen Beschreibung auch einige unterstützende Diagramme erwarten.

